% !TeX program = xelatex
\documentclass[aspectratio=169,11pt]{beamer}

\usepackage[UTF8,fontset=none]{ctex}
\IfFileExists{fonts/NotoSansSC-VF.ttf}{
  \setCJKmainfont[Path=fonts/,Extension=.ttf]{NotoSerifSC-VF}
  \setCJKsansfont[Path=fonts/,Extension=.ttf]{NotoSansSC-VF}
  \setCJKmonofont[Path=fonts/,Extension=.ttf]{NotoSansSC-VF}
}{
  \IfFileExists{C:/Windows/Fonts/NotoSansSC-VF.ttf}{
    \setCJKmainfont[Path=C:/Windows/Fonts/,Extension=.ttf]{NotoSerifSC-VF}
    \setCJKsansfont[Path=C:/Windows/Fonts/,Extension=.ttf]{NotoSansSC-VF}
    \setCJKmonofont[Path=C:/Windows/Fonts/,Extension=.ttf]{NotoSansSC-VF}
  }{
    \setCJKmainfont{FandolSong-Regular}
    \setCJKsansfont{FandolHei-Regular}
    \setCJKmonofont{FandolFang-Regular}
  }
}
\usepackage{amsmath,amssymb}
\usepackage{booktabs,tabularx,array}
\usepackage{tikz}
\usepackage{pgfplots}
\usetikzlibrary{arrows.meta,positioning}
\pgfplotsset{compat=1.18}

\definecolor{Ink}{HTML}{17324D}
\definecolor{Teal}{HTML}{147D7E}
\definecolor{Orange}{HTML}{D97732}
\definecolor{Red}{HTML}{B64C4C}
\definecolor{Green}{HTML}{4E8B57}
\definecolor{Mist}{HTML}{EDF3F6}
\definecolor{Warm}{HTML}{FFF2E8}
\definecolor{SoftGreen}{HTML}{EAF4EB}
\definecolor{GrayText}{HTML}{5B6570}

\usetheme{Madrid}
\usecolortheme{default}
\setbeamercolor{normal text}{fg=Ink,bg=white}
\setbeamercolor{structure}{fg=Teal}
\setbeamercolor{frametitle}{fg=white,bg=Ink}
\setbeamercolor{title}{fg=white}
\setbeamercolor{subtitle}{fg=white}
\setbeamercolor{author}{fg=white}
\setbeamercolor{institute}{fg=white}
\setbeamercolor{date}{fg=white}
\setbeamercolor{block title}{fg=white,bg=Teal}
\setbeamercolor{block body}{fg=Ink,bg=Mist}
\setbeamercolor{alertblock title}{fg=white,bg=Orange}
\setbeamercolor{alertblock body}{fg=Ink,bg=Warm}
\setbeamercolor{exampleblock title}{fg=white,bg=Green}
\setbeamercolor{exampleblock body}{fg=Ink,bg=SoftGreen}
\setbeamercolor{author in head/foot}{fg=GrayText,bg=white}
\setbeamercolor{date in head/foot}{fg=GrayText,bg=white}
\setbeamertemplate{navigation symbols}{}
\setbeamertemplate{headline}{}
\setbeamertemplate{itemize item}{\color{Teal}$\blacktriangleright$}
\setbeamertemplate{itemize subitem}{\color{Orange}$\bullet$}
\setbeamersize{text margin left=0.7cm,text margin right=0.7cm}
\setbeamertemplate{footline}{%
  \leavevmode%
  \hbox{%
    \begin{beamercolorbox}[wd=.78\paperwidth,ht=2.7ex,dp=1.2ex,leftskip=1em]{author in head/foot}%
      稀疏关系到个体化全局排序
    \end{beamercolorbox}%
    \begin{beamercolorbox}[wd=.22\paperwidth,ht=2.7ex,dp=1.2ex,rightskip=1em plus1fil]{date in head/foot}%
      \hfill\insertframenumber\,/\,\inserttotalframenumber
    \end{beamercolorbox}}
}

\title{CGR-v3.2：研究方法与已有结果}
\subtitle{稀疏局部关系学习与个体化全局排序}
\author{FSLR 项目汇报}
\institute{基于 Liu, Wang \& Luo (2026) 的行为机制复现}
\date{2026 年 7 月}

\begin{document}

{
\setbeamercolor{background canvas}{bg=Ink}
\begin{frame}[plain]
  \titlepage
\end{frame}
}

\section{研究方法与已有结果}

\begin{frame}{研究方法：汪康组}
  \small
  \begin{description}
    \item[实验任务]
      \begin{itemize}
        \item 每次完整实验（episode）包含 8 个物品（cue），物品与隐藏等级 A--H 的对应关系随机；
        \item 学习阶段呈现固定的 8 条局部关系，共重复 4 轮，因此有 32 次关系呈现；
        \item 每轮内的呈现顺序和左右位置随机，参与者只观察关系，不需要作答；
        \item 测验阶段比较全部 $\binom{8}{2}=28$ 个物品对，重复 10 轮，共作答 280 次，且没有反馈。
      \end{itemize}

    \item[状态更新范围]
      \begin{itemize}
        \item 每名虚拟被试开始时，模型都会重新初始化内部状态；
        \item 学习阶段每看到一条关系，就更新关系记忆和当前排序；
        \item 学习结束后固定最终排序；测验阶段只读取这一排序，不再更新。
      \end{itemize}
  \end{description}
\end{frame}

\begin{frame}{研究方法：汪康组}
  \small
  \begin{description}
    \item[模型输入]
      \begin{itemize}
        \item 学习阶段输入两个物品的身份、哪个物品等级更高，以及屏幕显示的带方向等级距离；
        \item 测验阶段只输入需要比较的两个物品；
        \item 模型看不到真实排名、正确答案、被试编号或测验反馈。
      \end{itemize}

    \item[关系记忆]
      每次看到一条关系时，模型以 $p_{\mathrm{enc}}=0.45$ 的概率将其记住。同一关系共呈现 4 次，因此至少成功记住一次的概率为
      \[
        1-(1-0.45)^4=0.9085.
      \]

    \item[个体先验]
      每个物品都有一个只在本次实验内有效的随机初始偏置（anchor）：
      $a_i\sim\mathcal N(0,\sigma_a^2)$，其中 $\sigma_a=0.30$。当学习证据不足时，这一偏置使不同虚拟被试形成不同但稳定的判断。

    \item[全局排序]
      模型用递归最小二乘（RLS）逐步整合关系信息，得到每个物品的等级分数；随后选择一个不违背已记忆关系的完整排序。所有测验题都从同一个最终排序中读取答案。
  \end{description}
\end{frame}

\begin{frame}{评估方法}
  \small
  \begin{description}
    \item[模拟对象]
      按论文的数据规模模拟 77 名虚拟被试。每名被试先观察 32 次局部关系，再完成 280 次无反馈判断。

    \item[准确率指标]
      分别统计总体正确率、8 个直接学习物品对、20 个未直接学习物品对，以及不同等级距离下的正确率。

    \item[结构指标]
      计算不同被试排序之间的 Kendall's $\tau$、同一被试的选择一致性、稳定错误占比和循环三元组数量。

    \item[分布指标]
      对每个物品对的跨被试准确率进行 Beta 分布拟合，记录按论文标准判定的双峰物品对数量；同时计算模型与人类 28 对准确率曲线的 Pearson 相关。
  \end{description}
\end{frame}

\begin{frame}{已有结果：总体表现}
  \begin{columns}[T,onlytextwidth]
    \begin{column}{0.5\textwidth}
      \centering
      \begin{tikzpicture}
        \begin{axis}[
          ybar,bar width=9pt,width=\textwidth,height=5.5cm,
          ymin=0.45,ymax=1.02,
          symbolic x coords={Overall,Learned,Unlearned,Tau},xtick=data,
          xticklabels={总体,已学习,未学习,个体间 $\tau$},
          x tick label style={font=\scriptsize},
          ylabel={指标值},
          legend style={at={(0.5,1.02)},anchor=south,legend columns=2,draw=none,font=\scriptsize},
          nodes near coords,nodes near coords style={font=\tiny,rotate=90,anchor=west}]
          \addplot[fill=Teal,draw=Teal] coordinates
            {(Overall,0.8524) (Learned,0.9140) (Unlearned,0.8277) (Tau,0.5807)};
          \addplot[fill=Orange,draw=Orange] coordinates
            {(Overall,0.8047) (Learned,0.9099) (Unlearned,0.7627) (Tau,0.5477)};
          \legend{人类,CGR-v3.2}
        \end{axis}
      \end{tikzpicture}
    \end{column}
    \begin{column}{0.5\textwidth}
      \begin{description}
        \item[直接学习]
          模型 $0.910$，人类 $0.914$，基本一致。
        \item[间接推断]
          模型 $0.763$，人类 $0.828$，模型偏低。
        \item[个体差异]
          模型 $\tau=0.548$，人类 $0.581$。
        \item[题目模式]
          模型与人类在 28 对题目上的准确率变化较相似，Pearson $r=0.779$。
      \end{description}
    \end{column}
  \end{columns}
\end{frame}

\begin{frame}{已有结果：距离效应}
  \begin{columns}[T,onlytextwidth]
    \begin{column}{0.5\textwidth}
      \centering
      \begin{tikzpicture}
        \begin{axis}[
          width=\textwidth,height=5.2cm,xmin=1,xmax=7,ymin=0.58,ymax=1.01,
          xtick={1,2,3,4,5,6,7},
          xlabel={客观等级距离},ylabel={正确率},
          grid=major,grid style={draw=Mist},
          legend style={at={(0.5,1.02)},anchor=south,legend columns=2,draw=none,font=\scriptsize}]
          \addplot[color=Teal,very thick,mark=*] coordinates
            {(1,0.7288)(2,0.8167)(3,0.8925)(4,0.9029)(5,0.9662)(6,0.9539)(7,0.9844)};
          \addplot[color=Orange,very thick,mark=square*] coordinates
            {(1,0.6338)(2,0.7643)(3,0.8673)(4,0.8948)(5,0.9242)(6,0.9396)(7,0.9429)};
          \legend{人类,CGR-v3.2}
        \end{axis}
      \end{tikzpicture}
    \end{column}
    \begin{column}{0.5\textwidth}
      \small
      \textcolor{Teal}{\textbf{共同模式}}\par
      人类和模型都表现出明显的距离效应：两个物品在真实排序中相距越远，判断越准确。

      \vspace{1.2em}
      \textcolor{Teal}{\textbf{主要差异}}\par
      当两个物品只相差 1--2 个等级时，模型的正确率低于人类，说明模型对相近物品的推断仍偏弱。
    \end{column}
  \end{columns}
\end{frame}

\begin{frame}{已有结果：个体化全局排序}
  \small
  \begin{description}
    \item[稳定性]
      模型的选择一致性为 $1.000$，人类为 $0.999$。同一虚拟被试在 10 轮测验中几乎始终作出相同选择。

    \item[稳定错误]
      模型的稳定错误占比为 $0.974$，人类为 $0.883$。这说明模型的错误通常来自固定的主观排序，而不是临时猜测；但模型比人类更少改变选择。

    \item[双峰分布]
      在 28 个物品对中，人类有 13 对、模型有 14 对呈现双峰分布。这表示不同被试会对同一物品对形成方向相反、但各自稳定的判断。

    \item[传递性]
      模型的所有答案都来自同一个完整排序，因此不会出现循环关系；人类数据中出现了 2 个循环三元组。
  \end{description}

  \begin{alertblock}{结果解释}
    模型同时表现出较高正确率、同一被试内部稳定、不同被试之间有差异，符合“个体化全局排序”的核心现象；但模型的选择比人类更固定。
  \end{alertblock}
\end{frame}

\begin{frame}{已有结果：模块消融}
  \footnotesize
  \begin{tabularx}{\textwidth}{p{0.22\textwidth}rrrrX}
    \toprule
    条件 & 总体 & 直接学习 & 个体间 $\tau$ & 双峰 & 主要结果 \\
    \midrule
    完整 CGR-v3.2 & 0.805 & 0.910 & 0.548 & 14 & 同时保留正确率、个体差异和稳定结构 \\
    去除个体先验 & 0.911 & 0.933 & 0.959 & 0 & 更接近客观排序，但不同被试趋同 \\
    去除关系记忆 & 0.491 & 0.474 & $-0.004$ & 28 & 缺少有效关系证据，结果退化 \\
    完整记忆 & 0.811 & 0.917 & 0.567 & 13 & 仅小幅提高准确率 \\
    去除共享排序读出 & 0.499 & 0.502 & 0.001 & 0 & 选择一致性降至 0.694，循环三元组升至 471 \\
    \bottomrule
  \end{tabularx}

  \vspace{0.8em}
  \begin{description}
    \item[个体先验]
      主要负责在相同学习输入下产生稳定的个体差异。
    \item[关系记忆]
      负责保存直接学习到的局部关系。
    \item[共享排序读出]
      负责被试内自洽性和选择的传递结构。
  \end{description}
\end{frame}

\begin{frame}{已有结果：探索路径与对齐}
  \footnotesize

  \centering
  \textbf{判断原则：行为相似不等于机制对齐；还要检验关键模块是否必要、个体差异来自哪里。}\par
  \vspace{0.35em}

  \begin{columns}[T,onlytextwidth]
    \begin{column}{0.48\textwidth}
      \begin{alertblock}{探索一：RetroModulRNN}
        \begin{itemize}
          \item \textbf{目标：}用神经调控可塑 RNN，让模型自行习得单个 episode 内的记忆与推断策略；
          \item \textbf{关键失败：}关闭可塑权重写入（plastic write）后，LPD 仅下降
                $0.0003$--$0.0004$；
          \item \textbf{判断：}模型表现几乎不依赖其核心快速可塑性，因此不能将该机制作为当前解释。
        \end{itemize}
      \end{alertblock}
    \end{column}
    \begin{column}{0.48\textwidth}
      \begin{block}{探索二：CGR-v1}
        \begin{itemize}
          \item \textbf{机制设想：}关系方向 $\rightarrow$ 逐步硬化全序；个体 anchor 用于补全欠定关系；
          \item \textbf{关键问题：}同 anchor、不同顺序的 $\tau=0.236$；不同 anchor、相同顺序的 $\tau=0.250$；
          \item \textbf{判断：}呈现顺序和早期硬化主导了差异，而非个体先验。结果相似，但来源没有对齐。
        \end{itemize}
      \end{block}
    \end{column}
  \end{columns}
\end{frame}

\end{document}
